% ============================================================
% front/resume.tex
% ============================================================
\chapter*{Résumé}
\addcontentsline{toc}{chapter}{Résumé}

AG Technologies, entreprise camerounaise spécialisée dans les solutions
digitales, exploitait un parc de serveurs physiques hébergeant des
applications critiques, piloté de façon manuelle et dispersée, sans
point de contrôle unique ni traçabilité homogène des actions effectuées.
Ce constat a motivé la conception d'AGT Infra, une plateforme interne
qui centralise la gestion de ce parc depuis un seul point d'entrée.

AGT Infra se positionne comme un chef d'orchestre plutôt que comme un
moteur : il pilote des briques déjà éprouvées (Docker, SSH, GitHub,
Nginx, Certbot) sans jamais les réimplémenter, en s'appuyant sur un
contrôle d'accès dynamique (RBAC/IBAC) où un refus explicite prévaut
toujours sur une autorisation plus large, ainsi que sur une traçabilité
systématique de toute action sensible. L'outil couvre la gestion des
serveurs et des conteneurs Docker, le déploiement continu à partir de
dépôts Git, la publication sécurisée d'applications sous HTTPS, la
supervision autonome par agent dédié, et une console web intégrée.

La majorité de ces domaines est aujourd'hui réalisée et vérifiée en
conditions réelles. Les sauvegardes automatiques, la montée en charge
automatique et un moteur de sécurité actif restent des besoins
identifiés mais non encore implémentés, et constituent les principaux
axes d'évolution du projet.

\bigskip
\noindent\textbf{Mots-clés :} infrastructure, DevOps, contrôle d'accès
dynamique (RBAC/IBAC), déploiement continu, conteneurisation,
supervision, orchestration SSH.